\documentclass[11pt,letterpaper]{article}
\usepackage[T1]{fontenc}
\usepackage{lmodern}
\usepackage[margin=1in]{geometry}
\usepackage{amsmath,amssymb,amsthm,mathtools}
\usepackage{microtype}
\usepackage{enumitem}
\usepackage{xcolor}
\usepackage[colorlinks=true,urlcolor=blue,linkcolor=black,citecolor=blue]{hyperref}
\hypersetup{pdftitle={Asymptotic improvement analysis for the October 2026 3 SAT preprint},pdfauthor={Research analysis},pdfsubject={Scoped barriers and a conditional improvement target}}
\pdfinfoomitdate=1
\pdftrailerid{}
\pdfsuppressptexinfo=15
\setlength{\parskip}{4pt}
\setlength{\parindent}{0pt}
\setlength{\emergencystretch}{2em}
\setlist{nosep,leftmargin=1.6em}
\newcommand{\poly}{\operatorname{poly}}
\newcommand{\F}{\mathbb F}
\newcommand{\Z}{\mathbb Z}
\newcommand{\bit}{\{0,1\}}
\newcommand{\pin}{\texttt{7ef5c0cc5ca4b0e072b7a050fca09ab24bc6fa3a}}
\newtheorem{proposition}{Proposition}
\newtheorem{lemma}{Lemma}
\newtheorem{hypothesis}{Unproved hypothesis}
\theoremstyle{remark}\newtheorem*{scope}{Scope}
\title{\vspace{-1.4em}Asymptotic improvement analysis\\for the October 2026 3 SAT preprint}
\author{Research report}
\date{27 October 2026}
\begin{document}
\maketitle
\vspace{-1.6em}
\begin{abstract}
\noindent
\textbf{No stronger unconditional worst-case bound is established.} The target is to beat the strongest Appendix D guarantee in OpenAI's \emph{The Exponential Time Hypothesis is False}, not to implement that guarantee in place of an older formal endpoint. Our principal conditional route would change the general-density exponent from $O(n\log\log n/\sqrt{\log n})$ to $O(n\sqrt{\log\log n/\log n})$, provided a degree-preserving compressed kernel and recursive closure can be constructed. They have not been constructed. A separate, stronger linear-wire compression hypothesis would remove the whole $\log\log n$ factor. We also prove restricted method barriers; none is a SAT lower bound. The fixed-density guarantee is not improved. Source inspection and finite tests do not verify the full preprint or its Lean dependencies.
\end{abstract}

\section{The target and the result}
Let $n$ be the number of declared Boolean variables and $m$ the number of input clauses. Width is fixed at three; the paper treats every fixed width $k$. Polynomial factors include reading the encoded input and exact bit arithmetic. Throughout this report, logarithms can be taken base two; a change of fixed base does not affect the displayed asymptotic comparisons. Small values of $n$ are handled separately.

The paper's strongest claimed general guarantee is
\begin{equation}\label{eq:baseline}
 T_{\rm paper}(n,m)=\poly(n,m)\,
 2^{O\left(n\log\log n/\sqrt{\log n}\right)}.
\end{equation}
For each \emph{fixed} density constant $D_0$, it already gives
\begin{equation}\label{eq:sparse}
 m\le D_0 n\quad\Longrightarrow\quad
 T_{\rm sparse}(n,m)=\poly(n,m)\,
 2^{O_{D_0}\left(n/\sqrt{\log n}\right)}.
\end{equation}
Both are exact-counting bounds in bit operations with polynomial space, and therefore also decide satisfiability \cite{paper}.

\textbf{Outcome.} The principal analyzed proposal is a recursive contraction interface with a global deformation degree $K$. If it existed for the sparse-query encoding, it would give the intermediate unrestricted-density exponent $O(n\sqrt{\log\log n/\log n})$, improving the general exponent by $\sqrt{\log\log n}$. The missing compressed kernel and closure properties are stated explicitly in Section~\ref{sec:conditional-degree-kernel}. A separate linear-wire compression hypothesis in Section~\ref{sec:conditional} would give the stronger exponent $O(n/\sqrt{\log n})$ for unrestricted density. Neither interface is proved. Neither proposal improves the paper's fixed-density exponent itself.

An upper bound $2^{O(f(n))}$ is a guarantee, not a lower bound on the old algorithm. Hence a sharper exponent would improve the certified asymptotic guarantee; it would not by itself establish that every concrete implementation runs faster on every input.

\section{Which source statement is being improved}
The inspected source is pinned to commit
\begin{center}\small\pin.\end{center}
The original preprint is dated 25 October 2026; its source title page states an earlier initial date and the same revision date. The repository describes its manuscripts as research outputs at differing verification stages \cite{repository}. The paper's title is not treated here as an independently established conclusion.

There are three distinct guarantees:
\begin{enumerate}
\item Theorem 1.1 and Appendix D: \eqref{eq:baseline}, and the stronger fixed-density bound \eqref{eq:sparse}.
\item The rational-interpolation argument in Sections 3--6: exponent
\[
 \rho(n)=n(\log\log n)^2/\sqrt{\log n}.
\]
\item The selected Lean endpoint \texttt{SqrtLogSpace.exact\_count\_and\_sat}, and the source declaration for its compiled-machine endpoint, refer to \texttt{SqrtLog.Statement.runtimeBound}. The selected canonical statement expands its exponential scale as
\[
 2^{A n (\log\log(n+16))^2/\sqrt{\log(n+16)}}.
\]
The machine and constants are quantified uniformly for a fixed width, rather than chosen separately for each instance \cite{leanfinal,leanstatement}.
\end{enumerate}
The third item is a source-level observation. No Lean kernel build or transitive axiom audit was performed for this report. In particular, implementing or formalizing Appendix D instead of this loglog-squared endpoint would reproduce an already stated paper bound. It would not satisfy the stronger target.

\section{Where Appendix D spends its exponent}
\subsection{The contraction interface}
The paper's Theorem D.1 contracts a bipartite Boolean tensor network with at most $Cs$ wires and maximum degree at most $\lceil s^{9/10}\rceil$. Its entries lie in
\[
 (\Z/2^{P_0}\Z)[V]/(V^\ell),\qquad 1\le P_0,\ell\le Cs.
\]
The number of vertices and description length are polynomially bounded in the applicable input parameters. For fixed constants, an entry must be evaluable in time
\[
 (s+\mathcal L)^b 2^{A s^{9/10}}
\]
and space $(s+\mathcal L)^b$, where $\mathcal L$ is the total graph and tensor-description length. Subject to these hypotheses, the theorem claims time
\begin{equation}\label{eq:contraction}
 (s+\mathcal L)^{O(1)}2^{O(s/\sqrt{\log s})}
\end{equation}
and polynomial space \cite{reduction}.

A direct clause-and-equality network uses $O(n+m)$ wires. This is enough for \eqref{eq:sparse}. At arbitrary density the sparse polynomial query construction instead uses $O(n\log\log n)$ wires. Substituting that size into \eqref{eq:contraction} produces \eqref{eq:baseline}. The extra factor is not merely a polynomial multiplication overhead.

\subsection{The local weighted enumeration}
At recursive depth $d$, the proof has $M=O(s8^{-d})$ wires and
\[
 u=\Theta((\log s)/2^d)
\]
retained valuation levels. At a node $t$, write $\delta=L/2^{v_2(t)}$ for the paper's mode-scale parameter. The number of variables of each positive weight $j$ is at most $6M/\delta$, and tapered interpolation yields a working precision
\[
 p_t=O(M\delta/u).
\]
Only mode monomials of total weight below $p_t$ can contribute. The resulting enumeration estimate is
\[
 \exp\!\left(O\!\left(\sqrt{(M/\delta)p_t}\right)\right)
 =2^{O(M/\sqrt u)}.
\]
Summing these exponents along a recursive path gives
\[
 \sum_{d\ge0} O\!\left(\frac{s8^{-d}}{\sqrt{(\log s)2^{-d}}}\right)
 =O(s/\sqrt{\log s}),
\]
since $\sqrt2/8<1$ \cite{contraction}.

The source uses immutable creation-precision tensor descriptions. Entrywise canonical lifts at later precision refer back to those fixed representatives. Discarding this discipline, or treating a computed residue as if it could be freely extended to higher precision, is not a valid optimization. Likewise, an entry evaluator is an actual algorithm with a cost bound, not a unit-cost oracle.

\section{A restricted obstruction to simpler enumeration bounds}
\begin{lemma}[Binary selections inside a weighted budget]\label{lem:binary}
Suppose an abstract mode inventory has $K$ distinct variables of each integer weight $1,\ldots,B$. If
\[
 K B(B+1)/2<P,
\]
then at least $2^{KB}$ different nonnegative exponent vectors have total weighted degree below $P$. All these vectors have ordinary degree at most $KB$.
\end{lemma}
\begin{proof}
For each of the $KB$ variables, independently select exponent zero or one. There are exactly $2^{KB}$ selections. The greatest possible weighted degree is the sum of all weights, namely $K\sum_{j=1}^B j=KB(B+1)/2$, and the greatest ordinary degree is $KB$.
\end{proof}

\begin{proposition}[The current numerical envelope permits exponential support]\label{prop:envelope}
Let $u\ge1$, $\delta\ge4\sqrt u$, and $M\ge2\delta$. Set
\[
 K=\lfloor M/\delta\rfloor,\qquad
 B=\lfloor\delta/(2\sqrt u)\rfloor,\qquad
 P=M\delta/u.
\]
An abstract inventory with $K$ variables of each weight $1,\ldots,B$ is consistent with the upper multiplicity bound $6M/\delta$. It has at least
\[
 2^{M/(8\sqrt u)}
\]
weighted-feasible exponent vectors of ordinary degree at most $M$.
\end{proposition}
\begin{proof}
Since $M/\delta\ge2$ and $\delta/(2\sqrt u)\ge2$, floor bounds give
\[
 \frac{M}{8\sqrt u}\le KB\le\frac{M}{2\sqrt u}\le M.
\]
Also
\begin{align*}
 \frac{KB(B+1)}2
 &\le \frac{M\delta}{8u}+\frac{M}{4\sqrt u}\\
 &\le\frac{M\delta}{8u}+\frac{M\delta}{16u}
 =\frac{3P}{16}<P.
\end{align*}
Apply Lemma~\ref{lem:binary}. The multiplicity condition holds because $K\le M/\delta$.
\end{proof}

\textbf{What this proves.} The multiplicity cap, weighted budget and ordinary-degree cap alone cannot imply a universally smaller exponent for the entire feasible-tuple universe. In particular, imposing the ordinary-degree constraint from a product of $M$ affine factors does not, by itself, remove the square-root-log scale. Sharper partition asymptotics or faster polynomial arithmetic alone do not establish a new log-power.

\textbf{What this does not prove.} The paper provides an upper inventory bound, not the required lower inventory on every actual network. The constructed envelope need not be realized with independent nonzero coefficients in a tensor instance. This is not a runtime lower bound for the paper's actual implementation, a lower bound for tensor contraction, or a lower bound for SAT. A new algorithm might exploit coefficient dependence, cancellation, or non-enumerative sharing that these numerical constraints ignore.

\section{A barrier for Hermite multiplicity retuning}
\label{sec:taper-barrier-audit}
This lemma concerns the normalization loss certificate in the paper's
Cauchy--Hermite construction. It is not a lower bound for SAT, tensor
contraction, or optimal Hermite recovery.

\begin{proposition}
Let $a\ge u\ge1$ be integers, $M>0$, $L=2^a$, and let
$q_0,\ldots,q_{u-1}$ be nonnegative integers.
Put $q_i=0$ for $u\le i\le a$. Suppose
\[
 L\sum_{i=0}^{u-1}2^{-i}q_i>M(2L-1).
\]
For $0\le i<u$, define
\[
 \ell_i=\sum_{j=1}^{a-i}\frac{j}{2^j}(q_i-q_{i+j}).
\]
Then some $i<u$ with $q_i>0$ satisfies
\[
 \ell_i>\frac{M}{4u}.
\]
Consequently, for any valuation-symmetric multiplicity profile satisfying
the degree-recovery condition, some retained node has normalization loss
$-v_2(E_t(0))=D_t\ell_i>D_tM/(4u)$.
\end{proposition}
\begin{proof}
Write $Q=\max_{0\le i<u}q_i$. Since $\sum_{i=0}^{u-1}2^{-i}<2$,
the degree condition gives $Q>M(1-1/(2L))\ge M/2$.
Let
\[
 W_i=\sum_{j=1}^{a-i}j2^{-j}\ge\frac12,\qquad
 d_i=\ell_i/W_i.
\]
Consider the following finite increasing Markov chain. From a state
$i<u$, jump to $i+j$, for $1\le j\le a-i$, with probability
$j2^{-j}/W_i$. Stop at the first state at least $u$. Every jump is
positive, so the stopping time $\tau$ is at most $u$, and
\[
 q_i-\mathbb E[q_{X_{t+1}}\mid X_t=i]=d_i.
\]
Start at a state with $q_i=Q$. Telescoping conditional expectations,
using $q_{X_\tau}=0$, yields
\[
 Q=\mathbb E\sum_{t=0}^{\tau-1}d_{X_t}
 \le u\max_{i<u}d_i.
\]
At a maximizing state, $d_i\ge Q/u>0$ and therefore
$\ell_i=W_id_i\ge Q/(2u)>M/(4u)$. If $q_i=0$, nonnegativity
of all later $q$ would give $\ell_i\le0$, so this state is retained.
The final assertion follows from the exact valuation identity in
the proof of Lemma D.5.
\end{proof}

The proposition does not require a monotone taper. It shows that replacing
the linear taper by another valuation-symmetric profile cannot improve the
\emph{worst normalization-loss bound} from order $MD/u$ to $o(MD/u)$.
It does not rule out a different recovery identity that exploits
cancellation among final Hermite coefficients, an asymmetric selection of
nodes, or a different deformation.

\section{Newton modes: exact formula, but no better logarithmic power}
For a crossing tile, put $h=2^{i+1}$ and
$A=T-y_0+x$, so $v_2(A)\le i$. For
$g(j)=(A-hj)^{-1}$, the forward-difference identity is
\[
 \Delta^b g(0)=
 \frac{b!\,h^b}{\prod_{r=0}^{b}(A-hr)}.
\]
It follows by induction from the definition of $\Delta$, or directly by
partial fractions. On a tile with integer indices $j=0,\ldots,d-1$,
Newton interpolation therefore represents the kernel exactly as
\[
 g(j)=\sum_{b=0}^{d-1}\binom jb\,\Delta^b g(0).
\]
The summands with $b>j$ vanish. This finite expression can replace the
geometric modes $j^b$ by the integral modes $\binom jb$.

Write $v=v_2(A)\le i$. Every denominator factor $A-hr$ has valuation
$v$, since $v_2(h)>v$. Thus
\[
 v_2(\Delta^b g(0))
 =-v+b(i+1-v)+v_2(b!).
\]
Multiplying by the Cauchy prefactor gives the certified weight
\[
 a+1-v+b(i+1-v)+v_2(b!).
\]
In the worst crossing valuation $v=i$, this becomes
\[
 a+1-i+b+v_2(b!)=a+1-i+2b-s_2(b),
\]
where $s_2(b)$ is the number of ones in the binary expansion of $b$.
Hence the weight remains linear in $b$. The factorial contributes only
a constant-factor improvement to the slope, not a superlinear weight.
The finite cap $b<d=D/2$ also does not remove modes at the partition
saddle scale $b\asymp D/\sqrt u$.

This is an exact alternative kernel representation worth investigating
for implementation constants. It does not prove an exponent better than
$O(M/\sqrt u)$, and therefore does not meet the stronger-than-Appendix-D
target.


\section{Decision shortcuts and the encoding obstacle}
Reducing the root modulus to two does not remove the current algorithm's internal precision cost. Its recurrence has the form
\[
 p_t=p+8\lceil3M/u\rceil\delta+8M(a+1).
\]
Even $p=1$ leaves the $\Omega(M\delta/u)$ term. Substituting the existing precision certificate into the mode estimate still gives the same $O(M/\sqrt u)$ exponent. This does not exclude a different parity-counting algorithm; it explains why parity alone does not improve this one.

There is also a precise limitation on lowering polynomial degree while preserving witnesses.
\begin{proposition}[Degree needed to preserve an odd solution count]\label{prop:odd}
For polynomials $P_1,\ldots,P_a$ over $\F_2$ on $n$ Boolean variables, if $\sum_i\deg P_i<n$, their number of common zeros is even. Consequently, a same-variable solution-set-preserving encoding of a uniquely satisfiable formula cannot always have total degree $o(n)$.
\end{proposition}
\begin{proof}
On the Boolean cube the common-zero indicator is
$Q(x)=\prod_i(1+P_i(x))$. Its degree is less than $n$. Every monomial in $Q$ omits some variable, so summing it over all Boolean assignments gives zero in $\F_2$. Thus $\sum_xQ(x)=0$. A unique common zero would instead give one.
\end{proof}
This statement does not constrain arbitrary decision reductions or reductions that intentionally change the witnesses.

The incidence encoding creates a different obstacle. If different rows use disjoint incidence variables and each has $r$ distinct cubic monomials, the product of $a$ row factors $(1+P_i)$ has at least $r^a$ distinct incidence monomials: independently choose one cubic term in each row. The rows' disjoint indeterminates prevent collisions or cancellation among these products. With $a=\Theta(n)$ and $r=\log^{\Theta(1)}n$, injectively indexing these monomials needs $\Omega(n\log\log n)$ bits.

This is not restricted to an arbitrary unrelated polynomial example. The bundled note \texttt{research\_decision/findings.md}, Section 2, gives a complete probabilistic construction from genuine 3-CNF violation polynomials. For $n\ge\lceil e^{10}\rceil$, choose $m=n\lceil\ln n\rceil$ clauses on distinct variable triples and $6n$ parity rows, including each clause independently with probability $4\ln(2m)/n$. A union bound over violation sets and assignments, together with row-size tail bounds, shows that some matrix preserves the exact solution set and has $\Theta(\log^2 n)$ distinct cubic terms per row. Its separated-incidence product has $2^{\Omega(n\log\log n)}$ nonzero monomials. The construction is an existence argument, not a proposed randomized solver or a description of the paper's deterministic greedy output.

Identifying incidence copies with their original Boolean variables may create very large cancellations and leaves at most $2^n$ multilinear supports. An algorithm that performs this quotient implicitly while retaining small local descriptions and low degree would evade the incidence-indexing obstruction. Such an algorithm has not been established here.

Other black-box retunings do not fix the problem. The supplied sparsification certificate, followed only by direct fixed-density contraction, optimizes for 3-SAT to exponent
\[
 O\!\left(\frac{n(\log\log n)^{2/3}}{(\log n)^{1/6}}\right),
\]
which is worse than \eqref{eq:baseline}. Branching on $t$ variables and paying the same black-box cost on $n-t$ variables gives exponent $t+Cf(n-t)$, where $f(n)=n\log\log n/\sqrt{\log n}$. Since $f'(n)\to0$, branching on $t\le n/2$ eventually increases that certificate; $t>n/2$ already costs exponentially more than $2^{Cf(n)}$. These conclusions concern those specific certificates, not all sparsification or branching algorithms.

% Research note; no unconditional improvement is asserted.
\section{A conditional degree-preserving compression theorem}
\label{sec:conditional-degree-kernel}
All logarithms in this section are base two, except in analytic partition estimates.
Let $s$ and $K$ be positive integers, with $s$ a size parameter, $\lambda=\lceil\log_2 s\rceil$, and suppose
$s/\lambda\le K\le s$. Constants multiplying these inequalities can be absorbed
by changing $s$ and the constants below. Put
\[
 H=\left\lceil\sqrt{sK/\lambda}\right\rceil,
 \quad u_d=\lfloor\varepsilon\lambda/2^d\rfloor,
 \quad a_d=2\sum_{j=0}^d u_j,
 \quad L_d=2^{a_d},\qquad \varepsilon=(128\cdot1024)^{-1}.
\]
The parameter $K$ is fixed globally; it need not decrease in a child.

\paragraph{Explicit hypothesis, not an established construction.}
Assume a class of encoded contraction objects admits the following uniform
recursive interface. Every object has $M\le Cs8^{-d}$ Boolean wires at depth $d$,
maximum degree at most $\lceil s^{9/10}\rceil$, and entries over the ring used in
Appendix D. Root precision and nilpotent truncation order are $O(s)$.
Here $\mathcal L$ denotes the total root description length. The graph and
immutable entry descriptions stored along any root-to-leaf path have total
length at most $(s+\mathcal L)^{O(1)}$, with a fixed uniform exponent, and all
entry evaluations, at their creation precisions, take
$\operatorname{poly}(s+\mathcal L)2^{O(s^{9/10})}$ time and polynomial space.
Each active object, $M>H$, must have all the following properties.
\begin{enumerate}
\item There is an exactly specified deformation $F(T)$ whose value at zero
is its contraction and whose degree is at most $K(2L_d-1)$. This is the key
new degree bound; the published bit-wire deformation only guarantees
$M(2L_d-1)$. For $S_d=\{-L_d,\ldots,L_d-1\}$ and
$\alpha_d=2^{a_d+1}$, all $F(t+\alpha_d z)$ are integral over the coefficient
ring. Their individual local factors and evaluations must be constructible
at polynomial bit cost, apart from the stated entry-evaluation cost.
\item At every retained node $t\ne0$ with $i=v_2(t)<u_d$, put
$D_t=L_d/2^i$. At any residue-field interpolation value of $z$, the evaluation
has an exact expansion modulo its working precision into mode monomials.
There are at most $cM/D_t$ modes of each positive integer weight $w$;
a mode monomial of weighted degree $b$ has coefficient divisible by $2^b$.
Modes and feasible exponent tuples have polynomial descriptions and can be
enumerated with polynomial delay and space. Merely showing $cN/D_t$ modes
for an uncompressed incidence count $N$ is not sufficient.
\item For each retained monomial, its contribution factors as the product
of a constant number of child contractions, each with at most $M/8$ wires.
The children remain in this class, with the same global degree parameter $K$,
maximum degree at most $\lceil s^{9/10}\rceil$, and the description and entry
cost bounds above. Entrywise canonical lifts must retain their creation-time
representatives, as in Appendix D. In particular, this condition prohibits
silently replacing a high-arity group tensor by an equality tree.
\item All local interpolation, packing, mode processing, rational recovery,
and description work per generated term is polynomial in $s+\mathcal L$;
the number of interpolation nodes, evaluation points, and groups per call
is polynomial in $s$ alone, independently of $\mathcal L$. Leaves can be summed over their Boolean wires directly, with the
stated entry costs. The coefficient-ring arithmetic uses polynomial-bit
coordinates and polynomial space at the precisions derived below.
\end{enumerate}

\paragraph{Conditional conclusion.}
Under this interface, contraction takes
\[
 \operatorname{poly}(s+\mathcal L)\,
 2^{O(\sqrt{sK/\log s})}
\]
bit operations and polynomial space. This is an implication from a new
recursive interface, not a proof that the interface exists for 3SAT.

\paragraph{Proof of the implication.}
Use the published tapered Hermite lemma with its degree parameter $M$
replaced by $K$. Thus
\[
 h=\lceil3K/u_d\rceil,\quad q_t=h\max(0,u_d-i),\quad
 p_t=p+8hD_t+8K(a_d+1).
\]
Here $K\ge u_d$ for all sufficiently large $s$. There are enough jets,
because $\sum q_t\ge L_dhu_d\ge3KL_d>K(2L_d-1)$; also $h\le4K/u_d$.
The proof of the source lemma then gives exactly the same integral recovery.
One may recover each jet by polynomially many residue-distinct evaluations
because the degree is $O(KL_d)=s^{O(1)}$.

Since $H\ge s/\lambda$, active depth is $O(\log\lambda)$; more precisely,
$8^d<Cs/H\le C\lambda$, so $u_d=\Omega_C(\lambda^{2/3})$.
Also $a_d\le4\varepsilon\lambda$, and the original packing inequalities
remain true: their right sides are $O(s^{0.902})=o(H)$ or smaller.
Summing previous precision increments gives
\[
 p\le s\lambda^{O(1)}2^{a_{d-1}}.
\]
For a retained node,
$D_t\ge2^{a_d-u_d}=2^{a_{d-1}+u_d}$.
As $K\ge s/\lambda$ and $2^{u_d}$ dominates every fixed power of $\lambda$,
the inherited precision and the $8K(a_d+1)$ term are absorbed, yielding
\[
 p_t=O(KD_t/u_d).
\]
All precisions remain polynomial in $s$ because $a_d=O(\log s)$.

The generating-function estimate from Appendix D, now with at most $cM/D_t$
variables of every weight, gives
\[
 \#\{\beta:\sum_Zw_Z\beta_Z<p_t\}
 \le \exp\bigl(O(\sqrt{Mp_t/D_t})\bigr)
 =\exp\bigl(O(\sqrt{MK/u_d})\bigr).
\]
At depth $d$ this exponent is at most
\[
 O\!\left(\sqrt{sK/\lambda}\,2^{-d}\right),
\]
because $M\le Cs8^{-d}$ and $u_d=\Theta_C(\lambda2^{-d})$.
Summing along paths, and accounting for the polynomial branching factors,
gives $\exp(O(\sqrt{sK/\lambda})+O(\lambda\log\lambda))$ calls.
The latter additive term is negligible since $\sqrt{sK/\lambda}\ge s/\lambda$.
Leaf work is bounded by $\operatorname{poly}(s+\mathcal L)
\exp(O(H)+O(s^{9/10}))$, with the same final exponent.
Depth-first accumulation and immutable polynomial descriptions give
polynomial space. This proves the conditional implication.

\paragraph{Consequent target for the sparse-query reduction.}
The source reduction has compressed wire budget $s=\Theta(n\nu)$ and retained
incidence degree $K=\Theta(n)$, where $\nu=\lceil\log_2\lceil\log_2 n\rceil\rceil$.
If its encoded objects and every recursive child admitted the above
interface, the source's other overheads would still be absorbed, giving
\[
 (n+m)^{O(1)}2^{O(n\sqrt{\log\log n/\log n})}
\]
for arbitrary-density exact counting and hence decision. This improves the
published exponent by a factor $\sqrt{\log\log n}$. It does not improve the
published fixed-density $O(n/\sqrt{\log n})$ exponent.
The retained incidence degree is not, by itself, a degree bound on the
multilinear extension of a binary subset index. No argument establishing
that implication or recursive closure has been found. Already the coefficient
table of the degree-one incidence polynomial $z_{2^b}$ on singleton indices,
encoded by $b$ bits in ordinary order, is the Boolean AND of all $b$ index
bits, of degree $b$. This elementary example rules out a generic degree
transfer based solely on the incidence degree.

\paragraph{A local candidate and its obstruction.}
For a categorical wire indexing $d$-element subsets of an $N$-element set,
consider $K_{S,T}(T_0)=\det W(T_0)[S,T]$, the $d$th exterior power of a base
matrix deformation. If $W(0)=I$, then $K_{S,T}(0)=\delta_{S,T}$; its degree is
at most $d\deg W$, and a block-diagonal $W$ induces block factorization after
group occupancies and determinant signs are tracked. Integrality of $W$
implies integrality of these determinants. Thus this is a genuine local
low-degree kernel, not yet the required recursive contraction algorithm.
Its ordinary Cauchy mode construction uses $N/D$, not
$\log\binom Nd/D$, modes per weight. Exterior antisymmetry can force a
rank-one mode to occur at most once in one determinant, but distinct-part
partition counts still have square-root entropy when the particle bound
is inactive. The missing compressed mode multiplicity and child-class
closure are substantive new lemmas.



\section{An exact sufficient compression lemma}\label{sec:conditional}
The following hypothesis isolates a sufficient missing theorem. It is deliberately stated with all interface costs rather than just a wire count.

\begin{hypothesis}[Uniform streaming compression]\label{hyp:compression}
Fix a function $W$ with $n\le W(n)\le n(\log n)^c$ for a fixed constant $c$. There is a deterministic polynomial-space procedure that, for every 3-CNF $F$, streams at most
\[
 \poly(n,m)2^{O(W(n)/\sqrt{\log W(n)})}
\]
Boolean tensor networks in the same total time. Uniform fixed constants guarantee the following for each network:
\begin{enumerate}
\item It is bipartite, has at most $CW(n)$ wires, at most polynomially many vertices in $W(n)$, and maximum degree at most $\lceil W(n)^{9/10}\rceil$.
\item Entries belong to $(\Z/2^{n+1}\Z)[V]/(V^\ell)$, where $1\le\ell\le CW(n)$. Write $\mathcal L$ for the total graph and evaluator-description length, which is polynomial in $n+m$.
\item Every entry at every Boolean port assignment is computable in time $(W(n)+\mathcal L)^b2^{AW(n)^{9/10}}$ and space $(W(n)+\mathcal L)^b$, with fixed constants $A,b$ independent of $F,n,m$. The immutable description may include the original input encoding; this adds no wires.
\item A specified polynomial-time linear functional of each contraction gives a residue; the sum of these residues is $\#F$ modulo $2^{n+1}$.
\end{enumerate}
\end{hypothesis}

\begin{proposition}[Conditional consequence]\label{prop:compression}
Assuming the paper's contraction theorem and Hypothesis~\ref{hyp:compression}, exact counting and decision take time
\[
 \poly(n,m)2^{O(W(n)/\sqrt{\log W(n)})}
\]
and polynomial space. In particular, $W(n)=O(n)$ would improve the unrestricted-density exponent to $O(n/\sqrt{\log n})$.
\end{proposition}
\begin{proof}
Apply the contraction theorem with $s=\lceil W(n)\rceil$ to one network at a time. Its initial precision $n+1$ is $O(s)$, and every other hypothesis is explicitly supplied. Multiplying the per-call bound by the number of networks changes only the fixed constant in the exponential rate. Apply the given functionals and accumulate modulo $2^{n+1}$. The true count lies between zero and $2^n$, so this residue determines the count exactly. Streaming preserves polynomial space.
\end{proof}

Any $W(n)=o(n\log\log n)$ in the stated range already improves the general exponent, because $\log W(n)=\Theta(\log n)$. The hard missing part is Hypothesis~\ref{hyp:compression}, not Proposition~\ref{prop:compression}. Merging the entire formula into one $n$-port tensor does not prove it: the degree cap fails, and evaluating that tensor after a purported degree reduction may conceal the original hard computation.

A concrete place to search is the paper's older original-variable interface before incidence expansion. It already uses $n$ Boolean variables and total input degree $O(n)$ but has the weaker loglog-squared analysis. A recursive replacement would need to prove that transformed factors retain controlled local degree, that shared measurement variables do not restore the incidence-code size, that local evaluation stays subdominant, and that interpolation precision remains controlled at every depth. None of these closure properties follows just from having $n$ original variables.



\section{What the audit and artifacts establish}
The audit inspected the claimed bounds and selected local proof mechanisms: residue-grid valuations, integral Hermite recovery, fixed representative lifting, recursive wire shrinkage, the degree invariant and depth-first space accounting. No immediate local algebraic contradiction was found in those checks. This is not a certification of every lemma, every reduction, every primitive operation, or the published ETH claim.

There is also a major practical qualification. Appendix D fixes
\[
 \varepsilon=1/(128\cdot1024)=2^{-17},\qquad
 u_0=\lfloor\varepsilon\lceil\log_2 s\rceil\rfloor.
\]
Thus $u_0=0$ when $\lceil\log_2 s\rceil<131072$. Even the displayed necessary root inequality $u_0\ge\lambda^{4/5}$ requires $\lambda\ge2^{85}$, before floors and deeper-call conditions are addressed. The proof handles a fixed exceptional range by direct summation. These constants reinforce that an asymptotic theorem is not evidence of a practical solver speedup.

The bundle includes the report source and PDF, pinned paper source files and PDF, selected Lean source excerpts, an independent audit, and original finite-check and rebuild scripts. The exact finite checks cover 27,990 taper profiles, 52,411 degree-admissible profile/degree cases, 388 rational normalizer products, 120 Newton difference identities, 420 interpolation values, 832 mode-envelope parameter cases, 313,600 explicit binary selections, and 1,312 affine polynomial systems. Separate exterior-power toy checks cover 36 identity entries, 288 integral substitutions, 216 block-locality comparisons, 65,536 shifted-mode subsets, and exact recurrence ratios. These local tests do not validate compressed mode multiplicities or recursive closure. They pass both normal Python and optimized Python; guard checks use explicit exceptions, not removable assertions. The archive contains no faster 3-SAT solver, no full implementation of Appendix D, and no new formal certificate. Their checks test only the stated finite identities, counts, source fingerprints and artifact integrity. They cannot prove an asymptotic theorem by experiment.

The report is compiled with shell escape disabled. Upstream Lean or repository programs are neither executed nor compiled. A fresh archive replay compares extracted PDF text and every rendered page against a clean report rebuild. Negative controls check rejection of changed, missing, extra, duplicate-path and path-traversal archive entries. A SHA-256 manifest provides internal consistency, not cryptographic authorship; the independently supplied hash of the sealed ZIP identifies this release. No dependency installation or network access is required by the replay scripts.

\section{The next proof obligation}
The useful next step is not to rename the existing Appendix D theorem as a new improvement. It is to prove the compressed-kernel interface and recursive closure in Section~\ref{sec:conditional-degree-kernel}, the stronger compression interface in Section~\ref{sec:conditional}, or to demonstrate a different method whose full bit-cost analysis is below \eqref{eq:baseline}. Any proposed proof must account for construction time, description size, local entry cost, precision growth, recursive invariants, and the surrounding sparse-query reduction. Until those obligations are met, the honest result remains a conditional research route plus limited negative results.

\begin{thebibliography}{9}\small\raggedright
\bibitem{paper} OpenAI. \emph{The Exponential Time Hypothesis is False}. Revision dated 25 October 2026. Theorem 1.1 and Appendix D. Pinned \href{https://github.com/openai/math/blob/7ef5c0cc5ca4b0e072b7a050fca09ab24bc6fa3a/preprints/The-Exponential-Time-Hypothesis-is-False-October-25-2026/paper.pdf}{paper PDF} and \href{https://github.com/openai/math/blob/7ef5c0cc5ca4b0e072b7a050fca09ab24bc6fa3a/preprints/The-Exponential-Time-Hypothesis-is-False-October-25-2026/build/main_full_version.tex}{main TeX source}.
\bibitem{reduction} OpenAI. Appendix D, reduction and contraction interface. Pinned \href{https://github.com/openai/math/blob/7ef5c0cc5ca4b0e072b7a050fca09ab24bc6fa3a/preprints/The-Exponential-Time-Hypothesis-is-False-October-25-2026/build/quantitative-reduction.tex}{\texttt{quantitative-reduction.tex}}.
\bibitem{contraction} OpenAI. Appendix D, proof of the contraction theorem. Pinned \href{https://github.com/openai/math/blob/7ef5c0cc5ca4b0e072b7a050fca09ab24bc6fa3a/preprints/The-Exponential-Time-Hypothesis-is-False-October-25-2026/build/quantitative-contraction.tex}{\texttt{quantitative-contraction.tex}}.
\bibitem{repository} OpenAI. Repository overview, including verification-stage qualifications. Pinned \href{https://github.com/openai/math/blob/7ef5c0cc5ca4b0e072b7a050fca09ab24bc6fa3a/README.md}{repository README}.
\bibitem{leanfinal} OpenAI. Selected formal endpoint source. Pinned \href{https://github.com/openai/math/blob/7ef5c0cc5ca4b0e072b7a050fca09ab24bc6fa3a/lean/OAI/Algorithms/SubexponentialSAT/SqrtLogSpace/Final.lean}{\texttt{SqrtLogSpace/Final.lean}} and \href{https://github.com/openai/math/blob/7ef5c0cc5ca4b0e072b7a050fca09ab24bc6fa3a/lean/OAI/Algorithms/SubexponentialSAT/SqrtLog/ExecutionNativeFinal.lean}{\texttt{SqrtLog/ExecutionNativeFinal.lean}}.
\bibitem{leanstatement} OpenAI. Explicit primitive-machine target expansion. Pinned \href{https://github.com/openai/math/blob/7ef5c0cc5ca4b0e072b7a050fca09ab24bc6fa3a/lean/OAI/Algorithms/SubexponentialSAT/SqrtLog/StatementCanonical.lean}{\texttt{SqrtLog/StatementCanonical.lean}}.
\end{thebibliography}
\end{document}
